\section{RegexMatcher~v2}\label{sec:design}

\subsection{From a Set of Regular Expressions to a Route Matcher}

RegexMatcher~v1 matches one string against a set of regular expressions in a single pass over a
prefix tree~\cite{stankov2024regex,regexmatcher2026v1}. The paper measures its engine at commit
\RegexMatcherVOneCommit. To route with it, a wrapper turns each pattern into a regular expression
with one capture group per parameter, and compiles one automaton per method. The lowest route
number among the matches wins, so precedence follows registration order. The wrapper finds a 405
by trying the other methods.

v2 adds a route matcher to the same library and leaves v1's engine unchanged. The route matcher is a
second set of headers in its own namespace, \texttt{matcher::route}, with its own build target. It
needs C++23 and the standard library only, and neither part of the library includes the other. Its
parser, table, lookup and compile-time table began as a pre-release that a development round
measured (Section~\ref{sec:prespec}). The paper measures v2 at commit \RegexMatcherCommit.

\subsection{Patterns and Their Meaning}\label{sec:semantics}

A pattern starts with a slash and splits at every slash into segments (Table~\ref{tab:segments}). A
literal segment may hold only the bytes that RFC~3986 allows in a path segment, the pchar
set~\cite{bernerslee2005rfc3986}. A pattern holds at most \VTwoMaxSegments{} segments and
\VTwoMaxParams{} parameters, and a parameter name appears once.

\begin{table}[H]
\caption{The segments of a pattern.\label{tab:segments}}
\centering\small
\begin{tabular}{@{}ll@{}}
\toprule
\textbf{Segment} & \textbf{Matches} \\
\midrule
\texttt{users} & the same bytes \\
\texttt{\{name\}} or \texttt{\{name:string\}} & one non-empty segment \\
\texttt{\{name:u64\}} & one segment that is an unsigned decimal number of type u64 \\
\texttt{\{name:i64\}} & one segment that is a signed decimal number of type i64 \\
\texttt{\{*name\}} & the non-empty rest of the path; only as the last segment \\
\bottomrule
\end{tabular}
\end{table}

Seven rules give a route its meaning. They are also the rules against which every arm's answers
are checked (Section~\ref{sec:checks}).
\begin{enumerate}
\item Paths are compared and captured as sent. Nothing is percent-decoded.
\item A parameter matches one segment of one or more bytes, so it accepts every pchar.
\item A catch-all matches a non-empty remainder of the path, slashes included, as sent.
\item A trailing slash is strict. The matcher never redirects.
\item At every segment, from left to right, a literal beats a typed parameter, which beats a plain
parameter, which beats a catch-all. The most specific route wins in any order of registration.
\item The matcher first looks for the request method's most specific route. Without one, it
answers ``method not allowed'' with the set of methods whose routes match, or ``not found''.
\item Two routes of one method that match exactly the same paths are an error. A compile-time
table fails to compile, and a build at run time returns an error.
\end{enumerate}
What HEAD and OPTIONS mean is the server's choice. The matcher answers for the method it is asked.

\subsection{The Table}\label{sec:table}

One table holds the routes of every method. A method whose routes are all literal is answered by an
exact-match table alone. It uses open addressing over a hash of the whole path and the method.
A slot holds the hash, the method, the route and the place of the path's bytes in a shared arena,
and a hit is confirmed by comparing those bytes.

Any other method gets a trie over path segments. Its fully literal routes stay in the trie too,
because the walk tries literal children first and so finds them first. A node records where its
literal children start and how many there are. It also records its typed, plain and catch-all
children, the route that ends at it, and how many kinds of child it has. An edge to a literal child
holds the first \KeyBytes{} bytes of the label as one machine word, the child, and the place of the
whole label in the arena. A node with up to \VTwoLinearEdges{} literal children scans them in order.
A wider node keeps them in a power-of-two array of slots, hashed by the word and the segment.

The four arrays of a table are the nodes, the edges, the exact-match slots and the arena. After the
build they lie in one block aligned to a page of \VTwoPageBytes{} bytes. Each array starts at the
offset that a table built at compile time gives it. So a table has the same layout in both forms,
and only its place in memory differs.

\subsection{The Lookup}\label{sec:lookup}

A lookup takes the table, the method and the path, and fills a result in place. The result holds
the status, the route, and up to \VTwoMaxParams{} captured values. Each value is a pointer and a
length into the request's path, so a lookup never allocates. On ``method not allowed'' the result
also holds the set of methods that match.

For a method whose routes are all literal, the lookup probes the exact-match table. Otherwise it
walks the method's trie one segment at a time. It finds the end of a segment by loading the path a
word at a time and searching the word for a slash, without reading outside the path. It then tries
the literal children by word and length, comparing the remaining bytes only for labels longer than
a word. Then it tests the typed children, then the plain parameter, then the catch-all. A node with
more than one kind of child leaves a choice point. If the walk fails deeper down, it returns to the
last choice and tries the next kind. A node with one kind of child never causes a return. Without a
route of the request's method, the lookup tries each other method that has routes and reports the
set that matches.

\subsection{Building a Table at Run Time}\label{sec:build}

The build takes a span of route specifications, each a method, a pattern and a route number. It
returns the table, or an error that names the failing route. The build parses
every pattern into one shared array of segment records, which holds only the segments that each
pattern has. A route's entry is then small: its method, route, specification, first segment and
segment count. The build sorts indices of entries rather than the entries themselves, and splits
them between the exact-match table and the tries. One scratch stack holds the literal children of
the node under construction. The arrays are reserved from counts known after the sort, and the
result is copied once into the page-aligned block.

\subsection{Checked When the Program Compiles}\label{sec:checked}

The checked front end passes a pattern as a template argument, a string literal held in a class
type~\cite{snyder2018p0732}. One static assertion per kind of error then stops the compilation with
a message that names the error, such as \texttt{route pattern: unclosed '\{'}. The front end derives
the C++ type of each parameter: \texttt{std::string\_view}, \texttt{std::uint64\_t} or
\texttt{std::int64\_t}. A route declaration pairs a method and a pattern with a handler. The
handler must take one argument per parameter, in order, with those types, and then the server's
context arguments. Two static assertions check the count and then the types. A generic lambda has
no fixed arity, so it is checked by whether it can be called with those arguments.

A table of declarations is built during constant evaluation by the same build function as at run
time. A malformed pattern or two routes of one method that match the same paths stop the
compilation. The compiler then names a function that tells the reason. Listing~\ref{lst:api} shows
the front end. The negative-compilation cases of H5(a) cover each of these errors.

\begin{lstlisting}[language=C++,caption={The checked front end. Each commented declaration at the
end, added to the array, stops the compilation with the reason shown.},label={lst:api},float=t]
#include <matcher/route.hpp>
using namespace matcher::route;

struct Ctx { /* the server's request and response */ };
std::string get_user(std::uint64_t id, Ctx& c);
std::string get_file(std::string_view name, std::string_view rest, Ctx& c);

using H = handlers<std::string, Ctx&>;
inline constexpr H::decl kApi[] = {
    H::get<"/users/{id:u64}", &get_user>(),
    H::get<"/files/{name}/{*rest}", &get_file>(),
    H::post<"/users/{id:u64}", &get_user>(),
};
inline constexpr auto kTable = make_route_table<kApi>();
static_assert(kTable.table.find(0u, "/files/a/b/c").route == 1);  // 0u is GET

// H::get<"/u/{id", &get_user>()        route pattern: unclosed '{'
// H::get<"/u/{id:u64}", &get_file>()    route handler: takes a different number
//                                        of parameters than the route pattern has
// H::get<"/users/{n:u64}", &get_user>() compile_time_table_has_duplicate_routes
\end{lstlisting}

A table built at compile time costs constant-evaluation steps and compiler memory. Clang and MSVC
bound the steps per evaluation~\cite{llvm2026clangmanual,microsoft2026constexpr}. The harness
compiles each table with a budget of \ConstexprSteps{} steps. Section~\ref{sec:hseven} reports
what each table needed.

\subsection{What the Rework Changed}\label{sec:rework}

Before the run, the engineering of v2 had three targets, each with a rule of adoption declared
before anything was measured. The first was the build time of large tables, which the development
round had found slower than the median competitor's. The rework of Section~\ref{sec:build} met
its rule, a lower build time in every development cell, and was adopted. It left the table's
contents identical, byte for byte, to those of the pre-release.

The second target was the instructions per lookup at \shape{mixed-disjoint} with \SizeTen{} routes,
where gin had run fewer. A node of gin's tree compares its whole prefix at once, which can span
several segments, while v2 takes one trie step per segment. The development round took this for the
cause. Chained edges, and two builds of a separate chain node, were built and measured on the lab
host. Each either raised the instructions of other cells or did not beat gin's count, which its rule
forbade, so none was adopted. The loss was therefore expected, and H2's text predicted it before the
run.

The third target was the gap between a table built at compile time and one built at run time. Both
arms of the harness now share one adapter, and both tables share one layout
(Section~\ref{sec:table}).

None of this changed the lookup's code. The out-of-line lookup that every v2 arm calls compiles to
identical machine code from the pre-release and from v2, with the run's compiler and flags. The
ablation arm of Section~\ref{sec:arms} keeps the pre-release's build and layout, so it measures what
the rework changed.
